\mxunit{expertes-c04}{Équations polynomiales}
\mxsubtitle{Second degré dans $\C$, factorisation, nombre de racines}

\begin{objectives}
Résoudre dans $\C$ une équation du second degré à coefficients réels ; factoriser un polynôme par $z-a$ lorsque $a$ est racine ; résoudre une équation du troisième degré dont une racine est connue ; savoir qu'un polynôme de degré $n$ a au plus $n$ racines.
\end{objectives}
\begin{prerequisites}
Second degré dans $\R$ (forme canonique, discriminant, somme et produit des racines) ; chapitres 1 et 2 ; somme des termes d'une suite géométrique.
\end{prerequisites}

\section{Automatismes}
\begin{fluency}[Sans calculatrice, cinq minutes]
\begin{enumerate}
\item Calculer le discriminant de $2x^2-3x+5$.
\item Factoriser $x^2-5x+6$.
\item Développer $(x-2)\left(x^2+2x+4\right)$.
\item Calculer $P(2)$ pour $P(x)=x^3-3x^2+4$.
\item Écrire $1+q+q^2+\dots+q^{n}$ sous forme d'un quotient, pour $q\neq1$.
\item Donner la somme et le produit des racines de $x^2-7x+10$.
\end{enumerate}
\end{fluency}
\begin{solution}[Automatismes]
\begin{enumerate}
\item $\Delta=9-40=-31$.
\item $(x-2)(x-3)$.
\item $x^3+2x^2+4x-2x^2-4x-8=x^3-8$.
\item $P(2)=8-12+4=0$.
\item $\dfrac{1-q^{n+1}}{1-q}$.
\item Somme $7$, produit $10$ (racines $2$ et $5$).
\end{enumerate}
\end{solution}

\mxnewpage
\section{Activité d'introduction}
\begin{activity}[Une racine en cache deux autres \mxchip{35 min}]
On considère $P(z)=z^3-2z^2+4z-8$.
\begin{enumerate}
\item Tracer à la calculatrice la courbe de $x\mapsto x^3-2x^2+4x-8$ sur $\intff{-3}{4}$. Combien l'équation $P(x)=0$ semble-t-elle avoir de solutions réelles ? En repérer une.
\item Vérifier que $P(2)=0$.
\item Déterminer des réels $a$, $b$, $c$ tels que, pour tout $z$, $P(z)=(z-2)\left(az^2+bz+c\right)$. On développera le membre de droite et on identifiera les coefficients.
\item En déduire toutes les solutions complexes de $P(z)=0$.
\item Calculer la somme et le produit des trois solutions. Comparer avec les coefficients de $P$.
\item Combien de solutions complexes un polynôme de degré $3$ peut-il avoir au maximum ? Formuler une conjecture pour le degré $n$.
\end{enumerate}
\end{activity}
\begin{solution}[Activité]
\begin{enumerate}
\item La courbe ne coupe l'axe des abscisses qu'une fois, en $x=2$ : une seule solution réelle.
\item $8-8+8-8=0$.
\item $(z-2)\left(az^2+bz+c\right)=az^3+(b-2a)z^2+(c-2b)z-2c$. L'identification donne $a=1$, $b-2a=-2$, $c-2b=4$, $-2c=-8$, soit $a=1$, $b=0$, $c=4$ (la dernière équation est cohérente). Donc $P(z)=(z-2)\left(z^2+4\right)$.
\item $z^2+4=(z-2i)(z+2i)$ : les solutions sont $2$, $2i$ et $-2i$.
\item Somme : $2$ ; produit : $2\times2i\times(-2i)=8$. Pour $z^3+az^2+bz+c$, on observe somme $=-a$ et produit $=-c$.
\item Au plus trois. Conjecture : au plus $n$ pour le degré $n$.
\end{enumerate}
\end{solution}
\begin{teachernote}[Conduite de l'activité]
L'identification des coefficients repose sur une propriété admise par le programme : une fonction polynôme nulle a tous ses coefficients nuls. La signaler explicitement. La question 5 prépare le problème (relations entre coefficients et racines).
\end{teachernote}

\mxnewpage
\section{Notes de cours}
\subsection{Second degré à coefficients réels}
\begin{theorem}[Résolution dans $\C$]
Soit $a$, $b$, $c$ trois réels avec $a\neq0$, et $\Delta=b^2-4ac$. L'équation $az^2+bz+c=0$ a :
\begin{enumerate}
\item si $\Delta>0$, deux solutions réelles $\dfrac{-b-\sqrt\Delta}{2a}$ et $\dfrac{-b+\sqrt\Delta}{2a}$ ;
\item si $\Delta=0$, une solution réelle $-\dfrac b{2a}$ ;
\item si $\Delta<0$, deux solutions complexes conjuguées $z_1=\dfrac{-b-i\sqrt{-\Delta}}{2a}$ et $z_2=\dfrac{-b+i\sqrt{-\Delta}}{2a}$.
\end{enumerate}
Dans tous les cas, $az^2+bz+c=a(z-z_1)(z-z_2)$, $z_1+z_2=-\dfrac ba$ et $z_1z_2=\dfrac ca$ (avec $z_1=z_2$ si $\Delta=0$).
\end{theorem}
\begin{proof}
La forme canonique $az^2+bz+c=a\left[\left(z+\frac b{2a}\right)^2-\frac\Delta{4a^2}\right]$ reste valable dans $\C$. Si $\Delta<0$, on écrit $\frac{\Delta}{4a^2}=\left(\frac{i\sqrt{-\Delta}}{2a}\right)^2$, puis on factorise la différence de deux carrés. Un produit de complexes est nul si et seulement si l'un des facteurs l'est (car tout complexe non nul est inversible). Somme et produit s'obtiennent en développant $a(z-z_1)(z-z_2)$ et en identifiant.
\end{proof}
\begin{example}[Résoudre $z^2-2z+5=0$]
$\Delta=4-20=-16=(4i)^2$, donc $z=\frac{2\pm4i}2$ : les solutions sont $1-2i$ et $1+2i$.
\end{example}

\subsection{Factorisation par $z-a$}
\begin{definition}[Fonction polynôme, racine]
Une \emph{fonction polynôme} de degré $n$ est une fonction $P:z\mapsto a_nz^n+\dots+a_1z+a_0$ définie sur $\C$, où $a_0,\dots,a_n$ sont des complexes et $a_n\neq0$. Un complexe $\alpha$ tel que $P(\alpha)=0$ est une \emph{racine} de $P$.
\end{definition}
\begin{keypoint}[Propriété admise]
Si une fonction polynôme est nulle pour tout complexe $z$, alors tous ses coefficients sont nuls. Par conséquent, deux fonctions polynômes égales ont les mêmes coefficients : on peut \emph{identifier} les coefficients.
\end{keypoint}
\begin{theorem}[Factorisation de $z^n-a^n$]
Pour tous complexes $z$, $a$ et tout entier $n\ge1$ :
\[z^n-a^n=(z-a)\left(z^{n-1}+az^{n-2}+\dots+a^{n-2}z+a^{n-1}\right)=(z-a)\sum_{k=0}^{n-1}a^kz^{n-1-k}.\]
\end{theorem}
\begin{proof}
On développe : $(z-a)\sum_{k=0}^{n-1}a^kz^{n-1-k}=\sum_{k=0}^{n-1}a^kz^{n-k}-\sum_{k=0}^{n-1}a^{k+1}z^{n-1-k}$. En posant $j=k+1$ dans la seconde somme, elle devient $\sum_{j=1}^{n}a^jz^{n-j}$. Les termes d'indice $1$ à $n-1$ se compensent ; il reste le terme $k=0$ de la première somme, $z^n$, et le terme $j=n$ de la seconde, $a^n$. On obtient $z^n-a^n$.
\end{proof}
\begin{theorem}[Factorisation par $z-\alpha$]
Soit $P$ une fonction polynôme de degré $n\ge1$ et $\alpha$ une racine de $P$. Il existe une fonction polynôme $Q$ de degré $n-1$ telle que, pour tout complexe $z$, $P(z)=(z-\alpha)Q(z)$.
\end{theorem}
\begin{proof}
Écrivons $P(z)=\sum_{k=0}^{n}a_kz^k$. Comme $P(\alpha)=0$,
\[P(z)=P(z)-P(\alpha)=\sum_{k=1}^{n}a_k\left(z^k-\alpha^k\right).\]
Par le théorème précédent, chaque $z^k-\alpha^k$ s'écrit $(z-\alpha)Q_k(z)$ avec $Q_k$ polynôme de degré $k-1$ de coefficient dominant $1$. Donc $P(z)=(z-\alpha)Q(z)$ avec $Q=\sum_{k=1}^na_kQ_k$, de degré $n-1$ et de coefficient dominant $a_n\neq0$.
\end{proof}
\begin{theorem}[Nombre de racines]
Une fonction polynôme de degré $n\ge1$ a au plus $n$ racines.
\end{theorem}
\begin{proof}
Par récurrence sur $n$. Une fonction polynôme de degré $1$, $a_1z+a_0$ avec $a_1\neq0$, a exactement une racine, $-\frac{a_0}{a_1}$. Supposons la propriété vraie au degré $n$, et soit $P$ de degré $n+1$. Si $P$ n'a pas de racine, c'est terminé. Sinon, soit $\alpha$ une racine : $P(z)=(z-\alpha)Q(z)$ avec $Q$ de degré $n$. Si $P(\beta)=0$, alors $\beta=\alpha$ ou $Q(\beta)=0$. Comme $Q$ a au plus $n$ racines, $P$ en a au plus $n+1$.
\end{proof}
\begin{theorem}[Racines conjuguées]
Si $P$ est à coefficients réels et si $\alpha$ est une racine de $P$, alors $\overline\alpha$ est aussi une racine de $P$.
\end{theorem}
\begin{proof}
Les coefficients étant réels, $P\left(\overline\alpha\right)=\sum a_k\overline\alpha^{\,k}=\overline{\sum a_k\alpha^k}=\overline{P(\alpha)}=0$.
\end{proof}
\begin{method}[Équation du troisième degré dont une racine est connue]
On factorise par $z-\alpha$ en identifiant les coefficients de $Q(z)=az^2+bz+c$, puis on résout $Q(z)=0$. Exemple : $z^3+8=0$ admet $-2$ pour racine et $z^3+8=(z+2)\left(z^2-2z+4\right)$. Le trinôme a pour discriminant $-12=\left(2i\sqrt3\right)^2$, d'où les solutions $-2$, $1-i\sqrt3$, $1+i\sqrt3$.
\end{method}
\begin{pitfall}
Le discriminant d'un trinôme à coefficients réels ne dit plus si l'équation a des solutions dans $\C$ : elle en a toujours. Il indique si elles sont réelles ou complexes conjuguées.
\end{pitfall}

\mxnewpage
\section{Exercices}
\begin{exercise}[Second degré]
Résoudre dans $\C$ : $z^2+9=0$ ; $z^2-4z+13=0$ ; $2z^2+2z+1=0$ ; $z^4-1=0$.
\end{exercise}
\begin{solution}
$z^2+9=(z-3i)(z+3i)$ : $\ens{-3i;3i}$. Ensuite $\Delta=-36$ : $2-3i$ et $2+3i$. Puis $\Delta=-4$ : $\frac{-2\pm2i}{4}$, soit $-\frac12-\frac12i$ et $-\frac12+\frac12i$. Enfin $z^4-1=\left(z^2-1\right)\left(z^2+1\right)$ : $\ens{1;-1;i;-i}$.
\end{solution}

\begin{exercise}[Se ramener au second degré]
\begin{enumerate}
\item Résoudre dans $\C$ l'équation $z^4+4z^2-5=0$ en posant $Z=z^2$.
\item Résoudre dans $\C$ l'équation $z+\dfrac1z=1$, pour $z\neq0$.
\end{enumerate}
\end{exercise}
\begin{solution}
\begin{enumerate}
\item $Z^2+4Z-5=(Z-1)(Z+5)$. Donc $z^2=1$ ou $z^2=-5$ : $\ens{1;-1;i\sqrt5;-i\sqrt5}$.
\item Pour $z\neq0$, l'équation équivaut à $z^2-z+1=0$, de discriminant $-3$ : $z=\frac{1-i\sqrt3}2$ ou $z=\frac{1+i\sqrt3}2$, tous deux non nuls.
\end{enumerate}
\end{solution}

\begin{exercise}[Une racine évidente]
On pose $P(z)=z^3-3z^2+4z-2$.
\begin{enumerate}
\item Vérifier que $1$ est racine de $P$.
\item Déterminer le polynôme $Q$ tel que $P(z)=(z-1)Q(z)$.
\item Résoudre $P(z)=0$ dans $\C$.
\end{enumerate}
\end{exercise}
\begin{solution}
\begin{enumerate}
\item $1-3+4-2=0$.
\item $(z-1)\left(z^2+bz+c\right)=z^3+(b-1)z^2+(c-b)z-c$, d'où $b=-2$ et $c=2$ : $Q(z)=z^2-2z+2$.
\item $Q$ a pour discriminant $-4$ : racines $1-i$ et $1+i$. Les solutions sont $1$, $1-i$, $1+i$.
\end{enumerate}
\end{solution}

\begin{exercise}[Une racine non réelle]
On pose $P(z)=z^3+z^2+4z+4$.
\begin{enumerate}
\item Vérifier que $2i$ est racine de $P$. En déduire une autre racine sans calcul.
\item Développer $(z-2i)(z+2i)$, puis factoriser $P(z)$ et résoudre $P(z)=0$.
\end{enumerate}
\end{exercise}
\begin{solution}
\begin{enumerate}
\item $P(2i)=-8i-4+8i+4=0$. Les coefficients sont réels, donc $-2i=\overline{2i}$ est aussi racine.
\item $(z-2i)(z+2i)=z^2+4$ et $P(z)=z^2(z+1)+4(z+1)=(z+1)\left(z^2+4\right)$. Les solutions sont $-1$, $2i$, $-2i$.
\end{enumerate}
\end{solution}

\begin{exercise}[Factoriser $z^3-27$]
\begin{enumerate}
\item Factoriser $z^3-27$ par $z-3$ à l'aide du cours.
\item Résoudre dans $\C$ l'équation $z^3=27$ et placer les solutions dans le plan complexe. Que remarque-t-on ?
\end{enumerate}
\end{exercise}
\begin{solution}
\begin{enumerate}
\item $z^3-3^3=(z-3)\left(z^2+3z+9\right)$.
\item Le trinôme a pour discriminant $9-36=-27=\left(3i\sqrt3\right)^2$ : $z=\frac{-3\pm3i\sqrt3}2$. Les trois solutions $3$, $\frac{-3-3i\sqrt3}2$, $\frac{-3+3i\sqrt3}2$ ont pour module $3$ et forment un triangle équilatéral inscrit dans le cercle de centre $O$ et de rayon $3$ (chapitre 5).
\end{enumerate}
\end{solution}

\begin{exercise}[Un paramètre]
Soit $m$ un réel. On considère l'équation $(E_m)$ : $z^2-2z+m=0$.
\begin{enumerate}
\item Selon les valeurs de $m$, préciser le nombre et la nature des solutions de $(E_m)$.
\item On suppose $m>1$. On note $A$ et $B$ les images des solutions. Déterminer $m$ pour que le triangle $OAB$ soit équilatéral.
\end{enumerate}
\end{exercise}
\begin{solution}
\begin{enumerate}
\item $\Delta=4-4m$. Si $m<1$ : deux solutions réelles $1\pm\sqrt{1-m}$ ; si $m=1$ : une solution, $1$ ; si $m>1$ : deux solutions conjuguées $1\pm i\sqrt{m-1}$.
\item $OA=OB=\sqrt{1+(m-1)}=\sqrt m$ et $AB=2\sqrt{m-1}$. Le triangle est équilatéral si et seulement si $m=4(m-1)$, soit $m=\frac43$ (qui vérifie $m>1$).
\end{enumerate}
\end{solution}

\begin{exercise}[Retrouver un polynôme]
Déterminer le polynôme $P$ de degré $3$, à coefficients réels et de coefficient dominant $1$, tel que $1+i$ soit racine de $P$ et $P(0)=4$.
\end{exercise}
\begin{solution}
Les coefficients étant réels, $1-i$ est aussi racine. On factorise $P$ par $z-(1+i)$ ; le quotient, de degré $2$, s'annule en $1-i$ (car $1-i\neq1+i$), donc on peut encore le factoriser par $z-(1-i)$. Il reste un facteur de degré $1$ et de coefficient dominant $1$ : $P(z)=(z-1-i)(z-1+i)(z-r)=\left(z^2-2z+2\right)(z-r)$ pour un complexe $r$. $P(0)=-2r=4$, donc $r=-2$ et $P(z)=\left(z^2-2z+2\right)(z+2)=z^3-2z+4$. Vérification : $P(0)=4$ et $P(1+i)=(-2+2i)-2-2i+4=0$.
\end{solution}

\mxnewpage
\section{Problème}
\begin{problem}[Coefficients et racines d'un polynôme de degré 3]
Soit $P(z)=z^3+az^2+bz+c$ un polynôme à coefficients complexes, admettant trois racines $z_1$, $z_2$, $z_3$ (non nécessairement distinctes), de sorte que $P(z)=(z-z_1)(z-z_2)(z-z_3)$ pour tout complexe $z$.
\begin{enumerate}
\item En développant, démontrer que
\[z_1+z_2+z_3=-a,\qquad z_1z_2+z_1z_3+z_2z_3=b,\qquad z_1z_2z_3=-c.\]
\item Vérifier ces relations sur le polynôme de l'activité, $z^3-2z^2+4z-8$.
\item Démontrer que $z_1^2+z_2^2+z_3^2=a^2-2b$.
\item On considère $R(z)=z^3-2z^2+3z-5$, dont on admet qu'il s'écrit $(z-z_1)(z-z_2)(z-z_3)$. Calculer $z_1^2+z_2^2+z_3^2$ sans résoudre l'équation.
\item En déduire que les racines de $R$ ne sont pas toutes réelles.
\end{enumerate}
\end{problem}
\begin{solution}[Problème]
\begin{enumerate}
\item $(z-z_1)(z-z_2)(z-z_3)=z^3-(z_1+z_2+z_3)z^2+(z_1z_2+z_1z_3+z_2z_3)z-z_1z_2z_3$. Par identification des coefficients (propriété admise), on obtient les trois relations.
\item Les racines sont $2$, $2i$, $-2i$ : somme $2=-(-2)$ ; somme des produits deux à deux $4i-4i+4=4$ ; produit $8=-(-8)$.
\item $(z_1+z_2+z_3)^2=z_1^2+z_2^2+z_3^2+2(z_1z_2+z_1z_3+z_2z_3)$, donc $z_1^2+z_2^2+z_3^2=a^2-2b$.
\item $a=-2$ et $b=3$ : $z_1^2+z_2^2+z_3^2=4-6=-2$.
\item Si les trois racines étaient réelles, la somme de leurs carrés serait positive. Elle vaut $-2$ : au moins une racine n'est pas réelle.
\end{enumerate}
\end{solution}
\begin{teachernote}[Statut du problème]
Les formules de Viète sont un « problème possible » du programme. Le fait que tout polynôme de degré $3$ se factorise ainsi dans $\C$ (théorème de d'Alembert-Gauss) est hors programme : il est admis dans l'énoncé. Pour $R$, la fonction réelle associée est strictement croissante (sa dérivée $3x^2-4x+3$ a un discriminant négatif) : une seule racine réelle, deux racines conjuguées non réelles.
\end{teachernote}
